% 03_VRL_PravdivostSlozenychVyroku.tex
% Priprava: 1. rocnik, hodina 03, Vyrokova logika.
% Navazuje na 02_VRL_CoJeVyrok.tex; po teto hodine nasleduji kvantifikatory.
% Metodicka inspirace: https://www.realisticky.cz/dil.php?id=17
% Relevantni lekce webu: 010403 Konjunkce a disjunkce, 010404 Implikace,
% 010405 Ekvivalence. Zadani a reseni jsou puvodni.
% Kompilace z Work/MAT: latexmk -cd -pdf -interaction=nonstopmode -halt-on-error 1_rocnik/pripravy/03_VRL_PravdivostSlozenychVyroku.tex
\documentclass[11pt,a4paper]{article}
% Spolecny styl pro ucitelske pripravy MAT (Work/MAT)
% Verze 1.1, 2026-10-23
\NeedsTeXFormat{LaTeX2e}
\ProvidesPackage{priprava}[2026/10/23 v1.1 Sablona ucitelskych priprav MAT]

% Zakladni typografie a cestina
\RequirePackage[T1]{fontenc}
\RequirePackage[utf8]{inputenc}
\RequirePackage[czech]{babel}
\RequirePackage{lmodern}
\RequirePackage{microtype}
\RequirePackage{amsmath,amssymb}
\RequirePackage{xcolor}
\RequirePackage{enumitem}
\RequirePackage[most]{tcolorbox}
\RequirePackage{tabularx,booktabs,array}
\RequirePackage{fancyhdr}
\RequirePackage{needspace}
\RequirePackage{geometry}
\RequirePackage[hidelinks]{hyperref}

\geometry{a4paper,top=17mm,bottom=18mm,left=23mm,right=23mm}
\definecolor{notegray}{gray}{0.39}
\setlength{\parindent}{0pt}
\setlength{\parskip}{2.5pt plus 1pt minus 0.5pt}
\setlist[itemize]{topsep=2pt,itemsep=1pt,parsep=0pt,leftmargin=1.45em}
\setlist[enumerate]{topsep=2pt,itemsep=2pt,parsep=0pt,leftmargin=1.65em}
\setcounter{secnumdepth}{0}
\renewcommand{\arraystretch}{1.12}

% Zahlavi: vlevo rocnik | globalni tema, vpravo cislo: nazev hodiny
% Zapati zustava zcela prazdne.
\pagestyle{fancy}
\fancyhf{}
\renewcommand{\headrulewidth}{0.35pt}
\renewcommand{\footrulewidth}{0pt}
\setlength{\headheight}{14pt}
\newcommand{\nastavPripravu}[4]{%
  \fancyhead[L]{\small #1\enspace\textbar\enspace #2}%
  \fancyhead[R]{\small #3: #4}%
}

% Minimalisticke cernobile boxy. Kazda definice zvlast.
\tcbset{commonlesson/.style={enhanced,breakable=false,colback=white,colframe=black,
  colbacktitle=white,coltitle=black,boxrule=0.45pt,
  arc=0.8mm,outer arc=0.8mm,left=3mm,right=3mm,top=1.6mm,bottom=1.6mm,
  toptitle=1mm,bottomtitle=0.7mm,fonttitle=\bfseries,
  attach boxed title to top left={xshift=2.5mm,yshift=-1.15mm},
  boxed title style={colback=white,colframe=white,boxrule=0pt,left=1mm,right=1mm}}}
\newtcolorbox{zadaniBox}[1][]{commonlesson,title={Zadání},#1}
\newtcolorbox{prikladBox}[1][]{commonlesson,title={Příklad},#1}
\newtcolorbox{definiceBox}[1][]{commonlesson,title={Definice},#1}

% Spolecny box pro zadani a jeho reseni, bez mezery mezi castmi.
% Pouziti:
%   \begin{prikladReseniBox}
%     Zadani ulohy.
%     \tcblower
%     \begin{align*} ... \end{align*}
%   \end{prikladReseniBox}
% Spodni cast ma jemny sedy podklad a delici linku; lze v ni sazet
% vice radku matematiky. Obe casti maji jen jeden spolecny ramecek.
\tcbset{solutionlesson/.style={
  commonlesson,
  enhanced jigsaw,
  breakable,
  colbacklower=black!3,
  segmentation style={solid,draw=black!25,line width=0.35pt},
  fontlower=\small\color{notegray},
  before lower={\noindent\textbf{Řešení:}\enspace},
  leftlower=3mm,rightlower=3mm
}}
\newtcolorbox{prikladReseniBox}[1][]{solutionlesson,title={Příklad},#1}
\newtcolorbox{zadaniReseniBox}[1][]{solutionlesson,title={Zadání},#1}

% Pomucky pro autora pripravy: sedive poznamky; \reseni pro starsi pripravy
\newcommand{\poznamka}[1]{\par\begingroup\small\color{notegray}\itshape
  \textnormal{\textbf{Poznámka pro učitele.}}\enspace #1\par\endgroup}
\newcommand{\reseni}[1]{\par\begingroup\small\color{notegray}
  \textbf{Očekávané řešení:}\enspace #1\par\endgroup}
\newcommand{\blok}[1]{\needspace{5\baselineskip}\subsection*{#1}}
\newcommand{\vyrok}[1]{\uv{#1}}

\nastavPripravu{1. ročník}{Výroková logika}{03}{Pravdivost složených výroků}

\begin{document}

\blok{1. Navázání na minulou hodinu}
Připomeň zápis $\neg p$, $p\land q$, $p\lor q$, $p\Rightarrow q$ a $p\Leftrightarrow q$.
Rozliš \emph{výrok} a jeho \emph{pravdivostní hodnotu}.

\begin{prikladReseniBox}
Nechť $p$: \vyrok{Číslo 8 je sudé.} a $q$: \vyrok{Číslo 8 je prvočíslo.}
Rozhodni o pravdivosti výroků $p\land q$, $p\lor q$ a $\neg q$.
\tcblower
Platí $p=1$, $q=0$. Proto
\[
 p\land q=0,\qquad p\lor q=1,\qquad \neg q=1.
\]
\end{prikladReseniBox}
\poznamka{Nejdřív vyžaduj zdůvodnění slovy. Pro disjunkci platí matematické \uv{nebo}: stačí aspoň jedna pravdivá část.}

\blok{2. Co určuje pravdivostní tabulka}
\begin{definiceBox}
\textbf{Pravdivostní tabulka} přehledně zachycuje pravdivostní hodnotu složeného výroku pro všechny možné kombinace pravdivostních hodnot jeho jednotlivých částí.
\end{definiceBox}

Pro dva výroky $p$, $q$ nastávají právě čtyři kombinace jejich pravdivostních hodnot:
\[
 (1,1),\qquad(1,0),\qquad(0,1),\qquad(0,0).
\]

\begin{center}
\small
\begin{tabular}{@{}cc@{\hspace{4mm}}ccccc@{}}
\toprule
$p$ & $q$ & $\neg p$ & $p\land q$ & $p\lor q$ & $p\Rightarrow q$ & $p\Leftrightarrow q$\\
\midrule
1 & 1 & 0 & 1 & 1 & 1 & 1\\
1 & 0 & 0 & 0 & 1 & 0 & 0\\
0 & 1 & 1 & 0 & 1 & 1 & 0\\
0 & 0 & 1 & 0 & 0 & 1 & 1\\
\bottomrule
\end{tabular}
\end{center}

\poznamka{Nechej žáky nejprve vyjmenovat všechny dvojice hodnot. Uspořádání řádků je jen systematická pomůcka, nikoli další matematické pravidlo.}

\blok{3. Kdy je složený výrok pravdivý?}
\begin{definiceBox}
\textbf{Negace} $\neg p$ má vždy opačnou pravdivostní hodnotu než výrok $p$.
\end{definiceBox}

\begin{definiceBox}
\textbf{Konjunkce} $p\land q$ je pravdivá právě tehdy, když jsou pravdivé \emph{oba} výroky.
\end{definiceBox}

\begin{definiceBox}
\textbf{Disjunkce} $p\lor q$ je pravdivá právě tehdy, když je pravdivý \emph{alespoň jeden} z výroků.
\end{definiceBox}

\begin{definiceBox}
\textbf{Implikace} $p\Rightarrow q$ je nepravdivá pouze tehdy, když $p$ je pravdivý a $q$ nepravdivý.
\end{definiceBox}

\begin{definiceBox}
\textbf{Ekvivalence} $p\Leftrightarrow q$ je pravdivá právě tehdy, když mají $p$ a $q$ stejnou pravdivostní hodnotu.
\end{definiceBox}

\poznamka{Největší zádrhel bývá implikace: při $p=0$ je $p\Rightarrow q$ pravdivá bez ohledu na hodnotu $q$. Nepleť logickou implikaci s příčinnou souvislostí.}

\begin{prikladReseniBox}
Posuď pravdivost věty \uv{Jestliže je 12 liché číslo, pak je 12 dělitelné třemi.}
\tcblower
Předpoklad $p$: \uv{12 je liché} má hodnotu $0$; závěr $q$: \uv{12 je dělitelné třemi} má hodnotu $1$. Implikace $p\Rightarrow q$ má tedy hodnotu $1$.
\end{prikladReseniBox}

\blok{4. Výpočet pravdivostní hodnoty formule}
\begin{prikladReseniBox}
Sestav pravdivostní tabulku výroku $A=(p\land\neg q)\lor q$.
\tcblower
Nejdříve spočítej $\neg q$, potom $p\land\neg q$ a nakonec disjunkci s $q$.
\begin{center}
\begin{tabular}{@{}ccccc@{}}
\toprule
$p$ & $q$ & $\neg q$ & $p\land\neg q$ & $A$\\
\midrule
1 & 1 & 0 & 0 & 1\\
1 & 0 & 1 & 1 & 1\\
0 & 1 & 0 & 0 & 1\\
0 & 0 & 1 & 0 & 0\\
\bottomrule
\end{tabular}
\end{center}
\end{prikladReseniBox}
\poznamka{Nech postupně doplňovat jednotlivé sloupce. Je užitečné zvýraznit poslední sloupec jako výsledek celé formule.}

\begin{zadaniReseniBox}
Doplň pravdivostní tabulku výroku $B=(p\lor q)\Rightarrow p$. Rozepiš nejprve hodnoty $p\lor q$.
\tcblower
\begin{center}
\begin{tabular}{@{}cccc@{}}
\toprule
$p$ & $q$ & $p\lor q$ & $B$\\
\midrule
1 & 1 & 1 & 1\\
1 & 0 & 1 & 1\\
0 & 1 & 1 & 0\\
0 & 0 & 0 & 1\\
\bottomrule
\end{tabular}
\end{center}
Implikace je nepravdivá jen ve třetím řádku ($1\Rightarrow0$).
\end{zadaniReseniBox}

\blok{5. Samostatné procvičení}
\begin{zadaniReseniBox}
\begin{enumerate}[label=\alph*)]
  \item Je-li $p=0$ a $q=1$, urči hodnoty $\neg p$, $p\Rightarrow q$ a $p\Leftrightarrow q$.
  \item Jsou-li $p=1$, $q=1$ a $r=0$, urči hodnotu $\neg p\lor(q\Rightarrow r)$.
  \item Najdi všechny dvojice hodnot $p$, $q$, pro které je $p\Leftrightarrow q$ pravdivé.
\end{enumerate}
\tcblower
\begin{enumerate}[label=\alph*)]
  \item $\neg p=1$, $p\Rightarrow q=1$, $p\Leftrightarrow q=0$.
  \item $\neg p=0$, $q\Rightarrow r=0$, tedy $0\lor0=0$.
  \item $(p,q)=(1,1)$ a $(p,q)=(0,0)$.
\end{enumerate}
\end{zadaniReseniBox}

\textbf{Shrnutí na tabuli:} Konjunkce: oba. Disjunkce: alespoň jeden. Implikace: nepravdivá jen $1\Rightarrow0$. Ekvivalence: stejné hodnoty.

\poznamka{Pokud zbude čas, nech sestavit tabulku také pro $\neg(p\land q)$. Na další hodině následují kvantifikované výroky.}

\end{document}
